\title{Direct Preference Optimization: Your Language Model is Secretly a Reward Model}

\begin{abstract}
Large-scale unsupervised language models (LMs) acquire extensive world knowledge and demonstrate some reasoning capabilities; however, fine-grained modulation of their outputs remains challenging because their training relies entirely on unsupervised objectives. Common strategies to enhance model steerability involve obtaining human annotations that rank the quality of model outputs and then adapting the unsupervised LM to reflect these preferences, typically employing reinforcement learning from human feedback (RLHF). Yet, RLHF poses practical difficulties, being both intricate and prone to instability, as it requires training a reward model to encapsulate human judgment, followed by reinforcement learning to adjust the original LM to optimize the modeled reward, all while restricting deviation from the initial model parameters. This work proposes a novel way to parameterize the reward model within the RLHF framework, which enables direct analytical extraction of the optimal policy corresponding to human preferences. This insight allows the standard RLHF problem to be reframed as an application of a straightforward classification loss. Our approach, termed \textit{Direct Preference Optimization} (DPO), is efficient, robust, and easy to use, dispensing with the need for sampling from the LM during fine-tuning or extensive hyperparameter adjustments. Through empirical evaluation, we demonstrate that DPO matches or surpasses existing techniques in aligning LMs with human preferences. In particular, DPO-based fine-tuning delivers stronger control over generated sentiment than PPO-based RLHF, and achieves equivalent or better performance in summarization and single-turn dialogue tasks, all while offering a simpler training and implementation process.
\end{abstract}

\section{Introduction}
Large-scale unsupervised language models (LMs) trained on massive corpora have demonstrated remarkable emergent properties~\citep{chowdhery2022palm, brown2020language, touvron2023llama, bubeck2023sparks}. Nonetheless, the training data for these models originates from humans with highly diverse intentions, levels of expertise, and objectives. As a result, some of these intentions and skillsets may be undesirable to replicate; for instance, although an AI-based programming assistant should \textit{recognize} common coding errors to provide corrections, its code generation ought to be influenced predominantly by the instances of exceptional programming found in its dataset, even if such instances are uncommon. In a similar vein, we might wish for a language model to be \textit{cognizant} of a misconception held by 50\% of individuals, but we would not want the model to propagate this error in half of its relevant outputs. Put differently, it is imperative to distinguish between the \emph{appropriate outputs and conduct} the model should exhibit and the much broader set of \textit{capabilities and knowledge} it has obtained. This distinction is central to developing AI that is reliable, high-performing, and amenable to user control~\citep{ouyang2022training}. Existing approaches predominantly employ reinforcement learning (RL) to adapt LMs toward human preferences. In this work, we demonstrate that the standard RL-based preference learning objective can be solved precisely by applying a straightforward binary cross-entropy loss, thereby streamlining the overall preference alignment process.

Generally, established techniques impart target behaviors to language models via carefully constructed datasets of human feedback that encode what users consider safe and effective. This process of aligning model preferences typically takes place after the initial broad unsupervised pre-training phase on a vast textual corpus. While the most direct technique for preference learning involves supervised fine-tuning on annotated high-quality responses, the most effective strategies to date leverage reinforcement learning from human or automated feedback (RLHF/RLAIF; \citep{christiano2017deep, bai2022constitutional}). RLHF algorithms first fit a reward function to a dataset annotated with human preferences, then perform RL to train the model to prefer outputs scoring highly with respect to this function—while also constraining the policy to avoid significant deviation from its original parameters. Although RLHF approaches yield models that excel at dialogue and code generation tasks, their implementation is considerably more involved than basic supervised learning, necessitating the training of several distinct language models and requiring in-loop sampling from model outputs, thereby incurring high computational demands.

This work demonstrates a method for directly tuning a language model in accordance with human preferences, eliminating the need for explicit reward modeling or reinforcement learning. We introduce \textit{{Direct Preference Optimization} (DPO)}, an algorithm that inherently optimizes the same reward maximization objective—constrained by KL-divergence—as typical RLHF methods, yet features a streamlined implementation and training process. Conceptually, DPO adjusts the model by enhancing the log-likelihood of favored responses relative to disfavored ones, while leveraging adaptive, sample-specific importance weights. These weights are crucial in preventing model collapse, an issue observed when simply using naive probability ratio-based updates. As with conventional approaches, DPO adopts a theoretical framework for modeling preferences (e.g., the Bradley-Terry model; \cite{bradley1952rankanalysis}) to evaluate the correspondence between a reward signal and human preference data. However, whereas prior methods employ the preference model to establish a loss for training a reward estimator, which then guides policy optimization, DPO applies a variable transformation that redefines the loss directly in terms of the policy parameters. Using a corpus of human-annotated preferences over model outputs, DPO enables policy optimization through a straightforward binary cross-entropy objective, effectively yielding the optimal policy under a latent reward function shaped by the preference annotations.

The primary contribution of this work is the presentation of Direct Preference Optimization (DPO), a straightforward and reinforcement-learning-free strategy for training language models using preference feedback. Empirically, our studies reveal that DPO performs on par with or better than established techniques—including PPO-based RLHF—across various preference-driven tasks, such as sentiment adjustment, summarization, and conversational dialogue, applied to language models containing up to 6B parameters.

Large-scale self-supervised language models have demonstrated the ability to tackle certain tasks without task-specific training (zero-shot) \citep{radford2019language}, or with limited task examples provided through prompt engineering (few-shot) \citep{gpt3,megatron,chowdhery2022palm}. Yet, substantial gains in downstream performance and adherence to user instructions are often achieved by fine-tuning these models on curated collections of instructions and human-produced responses \citep{mishra-etal-2022-cross,sanh2022multitask,chung2022scaling,thoppilan2022lamda}. This process, known as instruction-tuning, enables large language models (LLMs) to extrapolate to new instructions beyond those seen during tuning, thereby enhancing their flexibility and usefulness \citep{chung2022scaling}. 

Although instruction tuning has proven effective, collecting comparative human ratings of model outputs is typically less demanding than acquiring expert-level demonstrations. Consequently, subsequent research has shifted towards refining LLMs based on human preference datasets, resulting in improved outcomes across domains such as translation \citep{kreutzer-etal-2018-reliability}, summarization \citep{stiennon2022learning,ziegler2020finetuning}, narrative generation \citep{ziegler2020finetuning}, and instruction adherence \citep{ouyang2022training,ramamurthy2023is}. These approaches usually involve first fitting a neural reward model to align with the human preference data, employing preference modeling frameworks like the Bradley-Terry model \citep{bradley1952rankanalysis}. The language model is then further optimized to maximize the resulting reward through reinforcement learning techniques, predominantly REINFORCE \citep{williams1992reinforce}, proximal policy optimization (PPO; \cite{schulman2017proximal}), or related strategies \citep{ramamurthy2023is}.

A related stream of work builds on instruction-tuned LLMs, utilizing human feedback to automatically generate new synthetic preference data targeting specific criteria (e.g., safety, harmlessness) \citep{bai2022constitutional}. In these cases, the human involvement is minimal, limited to supplying guidelines—typically in the form of a rubric—to guide the LLM's annotations. 

Such methods embody the intersection of two research trajectories: one concerned with employing reinforcement learning to train language models for a wide spectrum of tasks~\citep{Ranzato2015SequenceLT,paulus2018a,wu2018learning}, and the other focused on techniques for learning from human judgments \citep{christiano2017deep,kupcsik2018learning}. However, despite the appeal of incorporating relative human preferences, directly optimizing large language models via reinforcement learning is still beset by substantial practical difficulties. This work addresses the challenge by introducing a theoretically grounded methodology for optimizing with relative preferences, circumventing the necessity for reinforcement learning.

The study of policy learning based on preferences, outside of linguistic contexts, has been explored within both the bandit and reinforcement learning paradigms, leading to a variety of methodological developments. When preferences or rankings over actions, rather than explicit rewards, are utilized for learning in contextual bandit problems, the resulting framework is termed the contextual dueling bandit (CDB; \cite{yue2012karmed,dudik2015contextual}). In scenarios where absolute reward signals are absent, theoretical analyses in CDBs replace the standard notion of optimality with the concept of a \textit{von Neumann winner}: a policy that ensures an expected victory rate of at least 50\% when compared pairwise with any other policy \citep{dudik2015contextual}. Despite this, the CDB framework typically assumes that preference feedback is delivered in an online fashion, whereas human preference learning generally leverages a static, offline dataset comprised of action pairs annotated with preference judgments \citep{yan2022human}. Relatedly, \textit{preference-based RL} (PbRL) explores the use of binary preferences generated from an opaque ‘scoring’ function, rather than standard reward signals \citep{BusaFekete2014,ruiz2023dueling}. A number of approaches for PbRL exist, with many first requiring the explicit estimation of this latent scoring (reward) function, often through modeling, followed by subsequent policy optimization \citep{jain2013learning,BusaFekete2014,christiano2017deep,sadigh2017active,kupcsik2018learning}, and some approaches support off-policy reuse of preference data. In contrast to these two-step procedures, we introduce a unified approach in which the policy is optimized in a single stage to directly accommodate observed preferences.

\section{Preliminaries}\label{section:prelims}

We revisit the RLHF pipeline as articulated by \citeauthor{ziegler2020finetuning}, as well as its extensions in later works \citep{stiennon2022learning, bai2022training, ouyang2022training}. The RLHF methodology is typically organized into three distinct stages: 1) supervised fine-tuning (SFT); 2) preference data collection coupled with reward modeling; and 3) reinforcement learning-based optimization.

\textbf{SFT}: The process commences with the supervised fine-tuning of a pre-trained language model on curated, high-quality datasets tailored to the target application (such as dialogue systems or summarization). This results in the fine-tuned model denoted as $\pi^\text{SFT}$.

\textbf{Reward Modelling Phase}: In the next phase, the SFT model is conditioned on prompts $x$ to generate pairs of candidate completions $(y_1, y_2)\sim \pi^\text{SFT}(y \mid x)$. Human annotators are then tasked with evaluating these pairs, expressing a preference by selecting the superior response; this feedback is represented as $y_w\succ y_l \mid x$, indicating $y_w$ is favored over $y_l$ for the prompt $x$. Here, $y_w$ and $y_l$ correspond to the chosen and non-chosen completions within each candidate pair. The generation of such preference data is assumed to be governed by an underlying, unobserved reward function $r^*(y, x)$. To model these preferences, several probabilistic frameworks have been utilized, with the Bradley-Terry (BT) model \cite{bradley1952rankanalysis} being particularly prevalent (although more expressive models, such as those from the Plackett-Luce family \citep{plackett1975analysis, luce2012individual}, may also be employed given sufficiently ranked data). According to the BT model, the distribution governing human preferences $p^*$ can be formally expressed as:

\begin{equation}\label{eq:bradley-terry}
    p^*(y_1\succ y_2 \mid x)=\frac{\exp\left(r^*(x, y_1)\right)}{\exp\left(r^*(x, y_1)\right) + \exp\left(r^*(x, y_2)\right)}.
\end{equation}

Suppose we have access to a fixed dataset of comparison tuples $\mathcal{D}=\bigl\{x^{(i)}, y_w^{(i)}, y_l^{(i)}\bigr\}_{i=1}^N$ drawn from the true preference distribution $p^*$. A reward model can then be expressed as $r_{\phi}(x, y)$ with learnable parameters $\phi$, which are trained via maximum likelihood estimation. This setup lends itself naturally to a binary classification framework, in which the associated objective becomes the negative log-likelihood loss:

\begin{equation}\label{eq:reward_model}
    \mathcal{L}_R(r_{\phi}, \mathcal{D}) = -\mathbb{E}_{(x, y_w, y_l)\sim \mathcal{D}}\bigl[\log \sigma(r_{\phi}(x, y_w)- r_{\phi}(x, y_l))\bigr]
\end{equation}

where $\sigma$ denotes the logistic sigmoid function. In large language models, $r_{\phi}(x, y)$ is commonly initialized using weights from the SFT model $\pi^\text{SFT}(y \mid x)$, to which a final linear head is attached, mapping the transformer's last layer output to a scalar representing the reward~\cite{ziegler2020finetuning}. Previous works mitigate high reward variance by normalizing rewards to enforce $\mathbb{E}_{x,y\sim \mathcal{D}}\left[r_\phi(x, y)\right] = 0$ for all $x$.

\textbf{RL Fine-Tuning Phase}: In the subsequent RL fine-tuning phase, the trained reward model supplies feedback to optimize the language model. Following the frameworks of \citet{jaques2017sequence, jaques2020human}, the learning objective can be written as

\begin{equation}\label{eq:RL}
\max_{\pi_{\theta}}  \mathbb{E}_{x\sim \mathcal{D}, y\sim \pi_{\theta}(y \mid x)}\bigl[r_{\phi}(x, y)\bigr] - \beta\mathbb{D}_{\textrm{KL}}\bigl[\pi_{\theta}(y\mid x)\mid \mid \pi_\text{ref}(y\mid x)\bigr],
\end{equation}

where $\beta$ regulates how far $\pi_\theta$ may diverge from a reference policy $\pi_\text{ref}$, generally set to the SFT model $\pi^\text{SFT}$. Typically, the initial parameters of $\pi_\theta$ are copied from $\pi^\text{SFT}$. The KL divergence constraint curbs the model from straying excessively from the input distribution the reward model can reliably evaluate, and also serves to promote output diversity, preventing over-concentration on a limited set of responses. Owing to the discrete sampling nature of text generation, direct gradient optimization is not applicable, necessitating reinforcement learning methods instead. Standard practice~\cite{ziegler2020finetuning, stiennon2022learning, bai2022training, ouyang2022training} is to use a reward of the form ${r(x, y) = r_{\phi}(x, y) -\beta (\log \pi_{\theta}(y\mid x) - \log \pi_\text{ref}(y\mid x))}$ and to employ PPO \cite{schulman2017proximal} for maximizing this objective.

\section{Direct Preference Optimization}\label{sec:DPO}

Addressing the limitations inherent in RL methods for large-scale model fine-tuning, our aim is to develop a direct approach for policy optimization grounded in preference data. Contrary to existing RLHF pipelines, which first fit a reward function and then tune the policy using reinforcement learning, our method introduces a reward parameterization for which the optimal policy can be determined analytically, circumventing the need for a reinforcement learning step. As we detail in the following, our central contribution lies in establishing an analytical correspondence from reward functions to optimal policies, enabling us to recast an objective originally defined over reward models into one directly over policies. By exploiting this change-of-variables, our framework forgoes learning an explicit reward model, while remaining compatible with established models

\textbf{Deriving the DPO objective.} To begin, we adopt the standard reinforcement learning objective as used in previous literature, as shown in Eq.~\ref{eq:RL}, with a general reward function $r$. Consistent with earlier studies~\citep{peters2007reinforcement, peng2019advantage, korbak2022reinforcement, go2023aligning}, it can be directly demonstrated that the optimizer of the reward maximization objective with a KL-divergence constraint, as given in Eq.~\ref{eq:RL}, yields:

\begin{equation}\label{eq:op_policy}
    \pi_r(y\mid x) = \frac{1}{Z(x)}\pi_\text{ref}(y\mid x)\exp\left(\frac{1}{\beta}r(x, y)\right),
\end{equation}

Here, the normalizing constant is defined as $Z(x) =\sum_{y}\pi_\text{ref}(y\mid x)\exp\left(\frac{1}{\beta}r(x, y)\right)$. A full proof is presented in Appendix~\ref{app:derivation1}. Notably, substituting the MLE-estimated reward $r_{\phi}$ in place of the true reward function $r^*$ does not eliminate the challenge of computing $Z(x)$~\citep{korbak2022reinforcement, go2023aligning}, as its estimation remains computationally intensive, limiting the practical usefulness of this formulation. Nevertheless, Eq.~\ref{eq:op_policy} may be algebraically rearranged to write the reward function in relation to its induced optimal policy $\pi_r$, the reference $\pi_\text{ref}$, and the normalization $Z(\cdot)$. Specifically, by applying the natural logarithm to both sides of Eq.~\ref{eq:op_policy} and reordering, we obtain:

\begin{equation}\label{eq:main_eq}
    r(x,y) =\beta \log \frac{\pi_r(y\mid x)}{\pi_\text{ref}(y\mid x)} + \beta \log Z(x).
\end{equation}

This change of variables can be applied to the true reward $r^*$ and its associated optimal policy $\pi^*$. Importantly, in the Bradley-Terry framework, outcomes are determined solely by reward differences between outputs, namely ${p^*(y_1 \succ y_2 \mid x) = \sigma(r^*(x, y_1) - r^*(x, y_2))}$. Replacing $r^*(x,y)$ using Eq.~\ref{eq:main_eq} within the preference equation (Eq.~\ref{eq:bradley-terry}) eliminates the partition function, so the resulting expression for human preference probability relies only on the optimal policy $\pi^*$ and the reference policy $\pi_\text{ref}$. As a result, the optimal RLHF policy $\pi^*$ in the Bradley-Terry model obeys:

\begin{equation}\label{eq:objective}
    p^*(y_1\succ y_2 \mid x)=\frac{1}{1 + \exp\left(\beta \log \frac{\pi^*(y_2\mid x)}{\pi_\text{ref}(y_2\mid x)} - \beta \log \frac{\pi^*(y_1\mid x)}{\pi_\text{ref}(y_1\mid x)}\right)}
\end{equation}

Appendix~\ref{app:derivation2} provides the full derivation. Although Eq.~\ref{eq:objective} employs the Bradley-Terry assumption, a similar strategy can yield equivalent forms for the broader class of Plackett-Luce models~\citep{plackett1975analysis, luce2012individual}, as discussed in Appendix~\ref{app:plackett_luce_models}.

Given that we now have a human preference likelihood expressed directly in terms of the policy rather than the reward, we are able to construct a maximum likelihood learning objective for a parameterized policy $\pi_\theta$. Mirroring the approach used in reward modeling (see Eq.~\ref{eq:reward_model}), this leads to the policy objective:

\begin{equation}\label{eq:optimum_model}
    \mathcal{L}_\text{DPO}(\pi_{\theta}; \pi_\text{ref}) = -\mathbb{E}_{(x, y_w, y_l)\sim \mathcal{D}}\left[\log \sigma \left(\beta \log \frac{\pi_{\theta}(y_w\mid

In this formulation, we approximate an implicit reward by adopting an alternative parameterization, wherein the optimal policy coincides with $\pi_\theta$. Notably, our method corresponds to fitting a reparameterized Bradley-Terry model, which ensures certain desirable theoretical guarantees, such as consistency under appropriate assumptions on the distribution of the preference data~\cite{bong2022generalized}. We elaborate further on the theoretical attributes of DPO and situate it within related literature in Section~\ref{sec:theory}.

\textbf{How does the DPO update operate?} To gain a mechanistic perspective on DPO, it is instructive to inspect the gradient of the loss function $\mathcal{L}_\text{DPO}$. Specifically, the gradient with respect to the parameters $\theta$ is expressed as follows:

\begin{multline*}\label{eq:gradient}
    \nabla_\theta \mathcal{L}_\text{DPO}(\pi_\theta;\pi_\text{ref}) = \\ -\beta\mathbb{E}_{(x, y_w, y_l) \sim \mathcal{D}} \bigg[\underbrace{\sigma(\hat{r}_\theta(x, y_l) - \hat{r}_\theta (x, y_w))}_\text{greater emphasis if reward ranking is inaccurate}\bigg[\underbrace{\nabla_\theta\log \pi(y_w \mid x)}_\text{boost likelihood of $y_w$} - \underbrace{\nabla_\theta\log\pi(y_l \mid x)}_\text{reduce likelihood of $y_l$}\bigg]\bigg],
\end{multline*}

Here, $\hat{r}_\theta(x, y) = \beta \log \frac{\pi_\theta(y \mid x)}{\pi_\text{ref}(y \mid x)}$ serves as the reward function implicitly induced by the language model $\pi_\theta$ in relation to the reference model $\pi_\text{ref}$ (further examined in Section~\ref{sec:theory}). Conceptually, the loss gradient acts to amplify the probability assigned to the preferred completions $y_w$ while diminishing the probability of the less preferred completions $y_l$. The impact of each training instance is modulated according to the degree to which the implicit reward $\hat{r}_\theta$ erroneously favors the dispreferred output, scaled by $\beta$, thereby reflecting both misranking severity and the KL constraint's tightness. Our empirical investigations underscore the critical role of this weighting scheme, as omitting the coefficient leads to model collapse, as documented in Appendix Table~\ref{tab:unlikelihood_generations}.

\textbf{Overview of DPO.} The typical DPO process involves: 1) Drawing completions $y_1, y_2 \sim \pi_\text{ref}(\cdot \mid x)$ for each prompt $x$, assigning human preference labels to generate an offline preference dataset $\mathcal{D} = \{x^{(i)}, y_w^{(i)}, y_l^{(i)}\}_{i=1}^N$; 2) training the language model $\pi_\theta$ by minimizing $\mathcal{L}_\text{DPO}$ with respect to the given $\pi_\text{ref}$, dataset $\mathcal{D}$, and target $\beta$ value. 
In real-world applications, preference datasets that are publicly available are generally reused, avoiding the need for fresh sample generation and collection of human preferences. Because such preference datasets are produced with $\pi^\text{SFT}$, we choose $\pi_\text{ref} = \pi^\text{SFT}$ wherever feasible. If $\pi^\text{SFT}$ is not at hand, we instead set $\pi_\text{ref}$ by fitting it to maximize the likelihood of the preferred completions ${(x, y_w)}$ in the data, specifically, ${\pi_\text{ref} = \argmax_{\pi}\mathbb{E}_{x, y_w \sim \mathcal{D}}\left[\log \pi(y_w \mid x)\right]}$. This approach serves to reduce the distribution mismatch between the actual reference distribution (which cannot be accessed) and the $\pi_\text{ref}$ used during DPO training. Implementation specifics and tuning choices are elaborated upon in Appendix~\ref{app:implementation}.

\section{Theoretical Analysis of DPO} This section offers a deeper theoretical perspective on the DPO method, clarifies its foundational basis, and connects the strengths of DPO to certain limitations observed in RLHF actor-critic approaches (including PPO~\cite{schulman2017proximal}).

\label{sec:theory}

\subsection{Your Language Model Is Secretly a Reward Model} DPO achieves both reward fitting and policy learning through a single maximum likelihood framework, dispensing with the need for explicit reward function learning or separate RL policy training. Specifically, the objective in Eq. \ref{eq:main_eq} aligns with the Bradley-Terry model, employing a reward structure $r^*(x, y) = \beta \log\frac{\pi^*_\theta(y \mid x)}{\pi_\text{ref}(y \mid x)}$. Training the parametric model $\pi_{\theta}$ can thus be seen as reward model optimization (see Eq. \ref{eq:reward_model}) after a suitable variable transformation. The remainder of this section develops the theoretical justification for this reparameterization, demonstrates that it places no restrictive constraints on the space of learnable reward models, and enables exact retrieval of the optimal policy. We start by establishing an equivalence among reward functions.

\begin{definition}
Two reward functions $r(x, y)$ and $r'(x, y)$ are defined as equivalent if and only if there exists a function $f$ such that ${r(x, y) - r'(x, y) = f(x)}$.     
\end{definition}

This equivalence criterion clearly constitutes an equivalence relation, effectively dividing the set of reward functions into equivalence classes. The following lemmas summarize key properties:

\begin{lemma}\label{lemma:same_prefrence}
Within the Plackett-Luce preference models, and specifically the Bradley-Terry framework, reward functions belonging to the same equivalence class yield identical preference distributions.
\end{lemma}

\begin{lemma}\label{lemma:same_policy}
Any two reward functions in the same equivalence class produce the same optimal policy in the constrained RL setting.
\end{lemma}

The arguments establishing these results are straightforward and can be found in Appendix \ref{app:lemma1}. The first lemma highlights the well-known under-specification in the Plackett-Luce model family \cite{plackett1975analysis}. Because of this ambiguity, identifiability constraints are generally introduced to ensure meaningful MLE estimates for the reward parameterization in Eq. \ref{eq:reward_model} \cite{bong2022generalized}. The second lemma asserts that all functions from a given equivalence class correspond to the same optimal policy. As such, for policy recovery, it suffices to reconstruct any representative reward from the target class. The following theorem, proven in Appendix~\ref{app:thm1}, formalizes a key characterization:

\begin{theorem}\label{thm:main}
Given mild regularity conditions, every equivalence class of reward functions compatible with the Plackett-Luce (and specifically Bradley-Terry) model structure admits a representation of the form ${r(x, y) = \beta \log \frac{\pi(y\mid x)}{\pi_\text{ref}(y\mid x)}}$ for some probability model $\pi(y\mid x)$ and a specified reference distribution $\pi_\text{ref}(y \mid x)$.
\end{theorem}

\begin{sproof}
Let $r(x, y)$ denote an arbitrary reward function, and suppose this determines an optimal policy $\pi_r(y \mid x)$ as characterized by Eq. \ref{eq:op_policy}. To see that a reward function from the equivalence class of $r$ can be written using the aforementioned reparameterization, define the transformation $f$ as  
\begin{equation}
    f(r; \pi_\text{ref}, \beta)(x, y) = r(x, y) - \beta\log\sum_{y}\pi_\text{ref}(y\mid x)\exp\left(\frac{1}{\beta}r(x, y)\right)
\end{equation}
This mapping $f$ acts to normalize the reward by subtracting the log partition function associated with $\pi_r$. As this normalization term depends only on $x$, the function $f(r; \pi_\text{ref}, \beta)(x, y)$ remains in the equivalence class defined by $r(x, y)$. Substituting $r$ with the right-hand side of Eq.~\ref{eq:main_eq} (which is valid for any $r$), we observe that $f(r; \pi_\text{ref}, \beta)(x, y) = \beta \log \frac{\pi_r(y\mid x)}{\pi_\text{ref}(y\mid x)}$. Thus, $f$ constructs a representative from the equivalence class of $r$ possessing the target parameterization, indicating that our reparameterized model incurs no loss of generality.
\end{sproof}

An alternative interpretation of Theorem~\ref{thm:main} is that it explicitly determines which representative of each equivalence class of reward functions the DPO reparameterization chooses—namely, the reward function fulfilling:

\begin{equation}\label{eq:lag_p}
     \sum_{y}\underbrace{\pi_\text{ref}(y\mid x)\exp\left(\frac{1}{\beta}r(x, y)\right)}_{=\pi(y\mid x)\text{, using Thm.~\ref{thm:main} reparam.}} = 1,
\end{equation}

meaning that $\pi(y\mid x)$ forms a valid probability distribution (all probabilities are nonnegative and collectively sum to 1). By examining Eq.~\ref{eq:op_policy}, it is apparent that Eq.~\ref{eq:lag_p} corresponds to the normalization constant, or partition function, of the optimal policy that arises from the reward function $r(x, y)$. The crucial insight behind the DPO method is its enforcement of constraints on the otherwise ambiguous Plackett-Luce (or, more specifically, Bradley-Terry) preference models. These constraints retain the expressiveness of the possible reward models while also ensuring that the optimal policy in Eq.~\ref{eq:op_policy} remains analytically manageable for any prompt $x$.

\subsection{Instability of Actor-Critic Algorithms} Our framework is also useful for identifying sources of instability in standard actor-critic methods such as PPO, which are widely applied in the RLHF context. Here, we concentrate on the RL fine-tuning phase described in Section \ref{section:prelims}, linking it to the control-as-inference approach~\cite{levine2018reinforcement} for the constrained RL setup presented in \ref{eq:RL}. Letting $\pi_{\theta}(y\mid x)$ denote the parameterized policy, the objective is to minimize $\mathbb{D}_{\text{KL}}[\pi_{\theta}(y|x) \mid \mid \pi^*(y\mid x)]$, with $\pi^*$ representing the optimal policy from Eq. \ref{eq:optimum_model} defined by reward $r_{\phi}(y, x)$. After appropriate mathematical manipulation, this leads to the following maximization objective:

\begin{equation}\label{eq:AC}
    \max_{\pi_{\theta}}\mathbb{E}_{\pi_{\theta}(y\mid x)}\left[\underbrace{r_{\phi}(x, y) -\beta\log\sum_{y}\pi_\text{ref}(y\mid x)\exp\left(\frac{1}{\beta}r_{\phi}(x, y)\right)}_{f(r_{\phi}, \pi_\text{ref}, \beta)} - \underbrace{\beta\log\frac{\pi_{\theta}(y\mid x)}{\pi_\text{ref}(y\mid x)}}_{\text{KL}}\right]
\end{equation}

This objective matches that targeted by earlier research \citep{ziegler2020finetuning, stiennon2022learning, bai2022training, ouyang2022training}, where the DPO-equivalent reward is applied for the reward class $r_{\phi}$. Within this context, the normalization component in $f(r_{\phi}, \pi_\text{ref}, \beta)$ can be seen as the soft value function corresponding to the reference policy $\pi_\text{ref}$. Although this normalization does not modify the optimum, omitting it can lead to high-variance policy gradients, destabilizing training. One approach to address the normalization term is to learn an explicit value function, but optimizing this can also be challenging. Alternatively, previous studies have normalized rewards using a baseline determined by human completions—essentially a Monte Carlo estimate based on a single sample for the normalization constant. The DPO formulation, however, produces a reward function that eliminates the need for such baselines.

\section{Experiments}
Here, we present empirical assessments of DPO's capability to learn policies directly from preference data. To begin, we consider a controlled text-generation environment to investigate DPO’s efficiency in balancing reward maximization and KL-divergence minimization from the reference policy, in comparison to widely used preference optimization methods like PPO. Subsequently, we apply DPO to larger-scale models and more complex RLHF challenges, encompassing summarization and dialogue. Our findings indicate that DPO, with minimal hyperparameter adjustment, typically achieves comparable or superior outcomes to competitive baselines such as RLHF with PPO and the method of selecting the best out of $N$ samples under a trained reward model. The experimental design is explained below, with supplementary methodological information available in Appendix~\ref{app:exp_details}.

\textbf{Tasks.} We investigate three distinct tasks involving open-ended text generation. Across all tasks, methods derive a policy using a preference dataset $\mathcal{D}=\bigl\{x^{(i)}, y_w^{(i)}, y_l^{(i)}\bigr\}_{i=1}^N$. The first task, \textbf{controlled sentiment generation}, utilizes $x$ as an initial segment from a movie review in the IMDb dataset \cite{maas-EtAl:2011:ACL-HLT2011}. Here, the goal is for the policy to output a continuation $y$ with an explicitly positive sentiment. For rigorous assessment, we create preference pairs over generated continuations, employing a pre-trained sentiment classifier to enforce $p(\text{positive}\mid x,y_w)>p(\text{positive}\mid x,y_l)$. The SFT baseline involves continued training of GPT-2-large until satisfactory performance is reached using training data from the IMDB dataset (see further explanation in App~\ref{app:sentiment_details}). In the second task, \textbf{summarization}, $x$ corresponds to a Reddit forum post; here, the aim is for the model to generate a summary $y$ capturing the post’s main arguments. In accordance with previous literature, we leverage the Reddit TL;DR summarization dataset \citep{volske-etal-2017-tl} and supplement with human preference annotations from \citeauthor{stiennon2022learning}. Our SFT model, adapted on human-written summaries of forum posts\footnote{\url{https://huggingface.co/CarperAI/openai_summarize_tldr_sft}}, is used with the TRLX \citep{leandro_von_werra_2023_7790115} system to facilitate RLHF. Notably, the preference data collected by \citeauthor{stiennon2022learning} were based on generations from a separate but similarly fine-tuned SFT model. The final task, \textbf{single-turn dialogue}, involves $x$ as a user query—ranging from technical science questions to personal advice. Here, the policy must generate a response $y$ that is informative and appropriate. We utilize the Anthropic Helpful and Harmless dialogue dataset \citep{bai2022training}, which features 170,000 dialogues between users and an automated agent. Each exchange concludes with two AI-generated responses and a human annotation marking the preferred answer. For this dataset, an existing SFT model is not available; instead, we construct our SFT model by continuing training of an out-of-the-box language model solely on the preferred responses.

\textbf{Evaluation.} We utilize two distinct strategies to assess our methods. For controlled sentiment generation, where the ground-truth reward function (a sentiment classifier) is accessible, we quantify each algorithm’s performance by charting the trade-off between attained reward and KL-divergence from the reference policy; this is feasible due to our knowledge of the true reward. Conversely, in real-world scenarios where the true reward function remains unknown, we gauge algorithmic performance based on their \textit{win rate} relative to a baseline, employing GPT-4 as a surrogate for human evaluators when measuring the quality of summaries and the helpfulness of single-turn dialogue responses. For summarization tasks, the baseline comprises the reference summaries from the test data, whereas in dialogue settings, the baseline is set by the test dataset’s preferred responses. Although prior research indicates that language models can surpass traditional metrics as automatic evaluators \citep{Chen2023ExploringTU}, we further substantiate our reliance on GPT-4 through a dedicated human study, presented in Sec.~\ref{sec:human-judgments}. Our findings reveal a high correspondence between GPT-4 and human ratings, often equaling or surpassing the level of agreement seen among human annotators.

\textbf{Methods.} Alongside DPO, we compare several established techniques for aligning language models with human preferences. For the summarization domain, we evaluate zero-shot prompting with \textbf{GPT-J} \citep{gpt-j}, and for dialogue, we employ 2-shot prompting using \textbf{Pythia-2.8B} \citep{biderman2023pythia}. We also include the \textbf{SFT} model and the \textbf{Preferred-FT} approach, in which models are fine-tuned through supervised learning on the selected completions $y_w$—sourced either from SFT (for controlled sentiment and summarization) or from a generic language model (in dialogue). Another technique, \textbf{Unlikelihood}~\citep{welleck2019neural}, directly encourages the model to increase the probability of $y_w$ while decreasing the likelihood assigned to $y_l$, introducing a tunable coefficient $\alpha\in[0,1]$ to weight the unlikelihood component. We further assess \textbf{PPO} \citep{schulman2017proximal}, which relies on a reward model trained from preference data, and \textbf{PPO-GT}, an oracle variant utilizing the actual reward in the controlled sentiment setup. For PPO-GT, we implement both a standard version \cite{leandro_von_werra_2023_7790115} and a refined version that applies reward normalization and optimizes hyperparameters for better results; these enhancements are likewise applied to PPO using learned rewards. Lastly, the \textbf{Best of $N$} strategy is examined, wherein $N$ candidate responses are sampled from the SFT model (or Preferred-FT in dialogue), with the highest-rewarded candidate—judged by a reward model learned from preferences—being selected. Although this approach often yields superior results, its reliance on sampling $N$ completions per test instance renders it computationally expensive for even moderate values of $N$.

\subsection{How well can DPO optimize the RLHF objective?}

Typical RLHF methods employ a KL-constrained reward maximization framework, which ensures that while the policy aims to maximize reward, it does not diverge significantly from the reference policy. Accordingly, a proper comparison of such algorithms should consider both the level of reward obtained and the associated KL divergence; superior reward accompanied by a considerable increase in KL is not always preferable. In Figure~\ref{fig:frontier-tldr-main}, we present the trade-off frontier between reward and KL across different algorithms in the sentiment analysis scenario. For each method, we conduct several training experiments, each time varying the policy's conservativeness parameter (target KL $\in\{3,6,9,12\}$ for PPO, $\beta \in \{0.05,0.1,1,5\}$, $\alpha\in\{0.05,0.1,0.5,1\}$ for unlikelihood, and random seed variations for preferred-FT), resulting in a total of 22 runs. At intervals of 100 training steps until the models converge, we assess the policies on a suite of test prompts, measuring both the average reward (using the true reward function) and the mean sequence-level KL divergence\footnote{Specifically, this refers to the cumulative sum of the KL divergences at each time step.} with respect to the reference policy, denoted as $\text{KL}\left(\pi\mid \mid \pi_\text{ref}\right)$. Our findings demonstrate that DPO defines the most effective frontier, obtaining superior rewards while keeping the KL low. This outcome stands out for several reasons. Firstly, DPO and PPO share the same optimization target, but DPO delivers significantly better efficiency—its reward-to-KL balance strictly outperforms that of PPO. Secondly, DPO's trade-off curve is superior even relative to PPO when PPO has direct access to ground-truth rewards (PPO-GT).

\subsection{Can DPO scale to real preference datasets?}
\label{sec:dpo-real-datasets}
We proceed to assess the fine-tuning capabilities of DPO on tasks involving summarization and single-turn dialogue. It is well-established that for summarization, automated metrics like ROUGE often display weak alignment with human preference~\citep{stiennon2022learning}, motivating prior efforts to fine-tune language models with PPO using human-labeled preferences to yield more effective summaries. To compare various approaches, we generate completions from the test partition of the TL;DR summarization dataset and determine the average win rate against reference completions from the same test set. Sampling is performed across a range of temperatures from 0.0 to 1.0 for all methods, with corresponding win rates illustrated in Figure~\ref{fig:frontier-tldr-main} (right). The same GPT-J SFT model\footnote{\url{https://huggingface.co/CarperAI/openai_summarize_tldr_sft}} is used as the base for DPO, PPO, and Preferred-FT fine-tuning. Results indicate that DPO attains an approximately 61\% win rate at a temperature setting of 0.0, outperforming PPO, which achieves about 57\% at its own optimal temperature (also 0.0). Furthermore, DPO reaches a higher peak win rate relative to the best-of-$N$ baseline. Notably, as DPO's $\beta$ hyperparameter was not exhaustively optimized in our experiments, these results may not reflect DPO's full capabilities. DPO also demonstrates considerably greater resilience to changes in sampling temperature when compared to PPO, whose performance can decline to that of the initial GPT-J model at elevated temperatures. Preferred-FT, meanwhile, yields minimal gains over the SFT model. In Section~\ref{sec:human-judgments}, we further conduct direct human evaluations comparing DPO and PPO: DPO samples produced at temperature 0.25 are chosen over PPO samples at temperature 0 in 58\% of pairwise judgments.

For single-turn dialogue experiments, we conduct evaluations on the portion of the Anthropic HH test split \citep{bai2022training} containing instances with a single round of human-assistant interaction. To assess method performance, GPT-4 is employed to compare the generated outputs against the test set’s preferred completions, from which we derive win rates for each approach. Given the lack of a standard SFT baseline for this benchmark, our pipeline begins with pre-trained Pythia-2.8B, leverages Preferred-FT to fit a reference model on the preferred completions—ensuring outputs are within-distribution—and subsequently trains using DPO. Additionally, we benchmark against two baselines: the top choice among 128 Preferred-FT completions (we observe that the Best of $N$ baseline saturates at $N=128$ for this task; see Appendix Figure~\ref{fig:best-of-n}) and a two-shot prompted variant of the original Pythia-2.8B model. Across all approaches, DPO matches or outperforms the strongest baselines for the optimal temperature settings per method. Furthermore, we assess an RLHF model, trained via PPO on Anthropic HH \footnote{\url{https://huggingface.co/reciprocate/ppo_hh_pythia-6B}} and sourced from a prominent repository \footnote{\url{https://github.com/CarperAI/trlx/tree/main/examples/hh}}, but find that neither prompt engineering nor temperature adjustment yields any advantage over the Pythia-2.8B baseline. Drawing on findings from TL;DR and the shared reward objective, we treat Best of 128 as an approximate stand-in for PPO’s capabilities. Overall, DPO emerges as the sole approach that is both efficient in computation and achieves an improvement over the Anthropic HH reference completions, offering performance on par with or exceeding the more resource-intensive Best of 128 baseline. Lastly, as illustrated in Figure~\ref{fig:dialogue-main}, DPO attains its optimal performance in a relatively brief number of steps.

\subsection{Generalization to a new input distribution}

To more rigorously assess how PPO and DPO perform when facing distributional shifts, we evaluate the policies trained on Reddit TL;DR summarization using data from a different domain: news articles from the test portion of the CNN/DailyMail dataset \citep{nallapati-etal-2016-abstractive}. For this evaluation, we use the optimal sampling temperatures found for TL;DR (0 and 0.25). Table~\ref{tab:ood} displays the outcomes. As before, we measure the win rate of each policy using GPT-4, comparing their outputs to the reference summaries from the dataset, adapting our GPT-4 (C) prompt for news articles by substituting “forum post” with “news article.” In this shifted domain, DPO maintains a strong advantage over PPO in terms of performance. These results offer early evidence that DPO-trained policies can generalize at least as robustly as those trained with PPO, despite DPO not leveraging the extra unlabeled Reddit TL;DR prompts utilized by PPO.

\subsection{Validating GPT-4 judgments with human judgments}
\label{sec:human-judgments}
We run a human subject study to determine the fidelity of GPT-4-based evaluations, focusing on the TL;DR summarization task and examining two different GPT-4 prompting strategies. The first prompt, \textbf{GPT-4 (S)} (simple), directs GPT-4 to identify the summary that better captures key information from the post. The second, \textbf{GPT-4 (C)} (concise), additionally asks which summary is more succinct; this prompt is considered because, with \textbf{GPT-4 (S)}, GPT-4 often favors longer and more redundant summaries compared to human annotators. (Refer to Appendix~\ref{app:prompts} for the full text of each prompt.) We design three separate pairwise evaluations—using the best-performing method (DPO, temp. 0.25), the lowest-performing (PPO, temp. 1.0), and an intermediate method (SFT, temp. 0.25)—to encompass a range of output qualities; each method is compared with the PPO model using its top sampling strategy (greedy). Analysis of these comparisons shows that, regardless of the prompt, GPT-4's preferences match those of human raters at a frequency comparable to human-human agreement rates, supporting the use of GPT-4 as a reliable stand-in for human evaluation (note: due to limitations in human reviewer availability, multiple judgments were only gathered for DPO and PPO-1 matchups). The \textbf{GPT-4 (C)} prompt, in particular, yields win rates that align more closely with human assessments, motivating its use for the primary results in Section~\ref{sec:dpo-real-datasets}. For further methodological information on the human evaluation—including the rater interface and the participant roster—see Appendix~\ref{app:human-study}.

\section{Discussion}
Preference-based learning represents an effective and scalable method for developing language models that are both competent and aligned with human values. In this work, we present DPO, a straightforward approach that enables language models to be trained from preferences without resorting to reinforcement learning methodologies. Instead of reshaping the preference learning problem to fit within a conventional RL framework simply to utilize existing RL solutions, DPO establishes a connection between language model policies and reward functions, allowing models to align with human preferences \textit{directly} through a basic cross-entropy objective, sidestepping both reinforcement learning and any loss in generality. Remarkably, DPO requires minimal hyperparameter adjustment and achieves performance that matches or surpasses traditional RLHF techniques, such as those based on PPO. This approach therefore substantially lowers the practical barriers for training language models using human preferences.

\textbf{Limitations \& Future Work.} Our findings open up several avenues for further investigation. One open question concerns the out-of-distribution generalization capabilities of DPO-trained policies relative to approaches that utilize explicit reward functions. Preliminary evidence indicates that DPO models generalize comparably to PPO-based models, yet systematic evaluation remains to be done. Additionally, it remains to be seen whether self-labeling strategies employing the DPO policy can leverage unlabeled data as effectively. Another area of inquiry is the behavior of reward over-optimization in the context of direct preference optimization, particularly in relation to the modest drop in performance observed in Figure~\ref{fig:dialogue-main}-right and whether this is attributable to such phenomena. While our analysis covers models up to 6B parameters, future research may examine the feasibility and implications of extending DPO to substantially larger, state-of-the-art models. With respect to evaluation practices, our results reveal that GPT-4-derived win rates are sensitive to the evaluation prompt, suggesting the need for research into optimizing automated system judgment elicitation. Finally, DPO has a wide range of potential applications extending beyond language model training from preferences, including the training of generative models across various domains.